A distribution over coherent repairs is not generally equivalent to choosing a new repair hypothesis at every token. When two diagnoses support mutually incompatible edits, token-wise mixing can produce continuations supported by neither diagnosis. We propose Particle-Belief Program Filtering (PBPF) to connect sequential uncertainty estimation with coherent program repair. For each candidate version, PBPF represents a learned latent failure posterior with weighted particles, updates it using ordered visible execution outcomes, and applies transition and proposal corrections only when a patch creates a new version. One posterior component conditions every token of a repair through a learned soft prefix attached to a frozen actor. Training likewise uses a mixture of complete sequence probabilities, rather than a product of token-level mixtures. We specify an evaluation that separates the validity of the belief from the utility of its induced action: an exact finite audit, calibrated prediction of evaluator-withheld outcomes, and equal-budget hidden-test repair. MAP, posterior-mean, token-remix, and full-ensemble controls are designed to test the proposed coherence mechanism alongside matched deterministic memories. All empirical outcomes remain pending, and the manuscript makes no claim that multimodality or sample-once conditioning has already improved repair. The contribution is a concrete probabilistic formulation and a falsifiable protocol for determining when preserving diagnostic uncertainty matters for code generation.
